% Compile with: uplatex file.tex && dvipdfmx file.dvi
%           or: latexmk -uplatex file.tex
\documentclass[dvipdfmx]{article}
\usepackage{amsmath, amssymb}
\usepackage{tikz}

\title{scenario\_rotating\_frame}
\date{\today}

\begin{document}
\maketitle

Register the lab-frame spin Hamiltonian $H_{\rm lab}=-\omega_0S^z_1-\omega_1S^x_1$ ($\omega_0$: the Zeeman frequency, $\omega_1$: the strength of the rotating field).

$H_{lab} = -\omega_0 S^z_1 - \omega_1 S^x_1$

\begin{equation}
  H_{lab} = -\omega_0 S^{z}_{1} - \omega_1 S^{x}_{1}
  \label{eq:E1}
\end{equation}

To move into a frame rotating at angular frequency $\omega$, add the frame term $+\omega S^z$ to the lab-frame Hamiltonian, giving the rotating-frame Hamiltonian $H_{\rm rot}$.

$H_{rot} = E1 + \omega S^z_1$

\begin{equation}
  H_{rot} = \left( \omega - \omega_0 \right) S^{z}_{1} - \omega_1 S^{x}_{1}
  \label{eq:E2}
\end{equation}

From the coefficients (field components) of $H_{\rm rot}=-\mathbf B\cdot\mathbf S$, directly find the magnitude of the effective field felt by the spin, $\omega_{\rm eff}=\sqrt{(\omega-\omega_0)^2+\omega_1^2}$.

Find the magnitude of the effective field of E2

\begin{equation}
  \omega_{\mathrm{eff}} = \sqrt{ \left(\omega - \omega_0\right)^2 + \omega_1^2 }
  \label{eq:E3}
\end{equation}

Compute the commutator of the rotating-frame Hamiltonian $H_{\rm rot}$ and $S^z$ to obtain the equation of motion for precession in the rotating frame.

Find the commutator of E2 and $S^z_1$

\begin{equation}
  [H_{rot}, S^z_1] = i \omega_1 S^{y}_{1}
  \label{eq:E4}
\end{equation}

\end{document}